%% kapa: exam=1 round=27 year=2016 period=1 form=A subject=economics subject_ko=경제학원론 seq=16 number=56
%% source: Q-Net 감정평가사 기출문제 (cst00309) artlSeq=5169402 file=extracted/1008/5169402/제27회 감정평가사 1차 문제/제27회 감정평가사 1교시문제(A형).pdf
%% conversion: kordoc 4.24.51 -> tools/gen1.py (auto); review=proofread-sonnet; flags=none
\begin{question}{56}
다음은 A국과 B국의 교역관계에 대한 보수행렬(payoff matrix)이다. 이에 관한 설명으로 옳은 것은? (단, 보수쌍에서 왼쪽은 A국의 보수이고, 오른쪽은 B국의 보수이다.)
\par{\small\centering
\begin{tabular}{|c|c|c|c|}\hline
\multicolumn{2}{|c|}{} & \multicolumn{2}{c|}{B국} \\ \cline{3-4}
 & 전략 & \makebox[3cm]{저관세} & \makebox[3cm]{고관세} \\ \hline
\multirow{2}{*}{A국} & 저관세 & (250, 250) & (300, 100) \\ \cline{2-4}
 & 고관세 & (100, 300) & (200, 200) \\ \hline
\end{tabular}\par}
\begin{choices}
  \item 내쉬균형은 2개이다.
  \item 내쉬균형에 해당하는 보수쌍은 (200, 200)이다.
  \item 우월전략균형에 해당하는 보수쌍은 (100, 300)이다.
  \item A국의 우월전략은 고관세이다.
  \item B국의 우월전략은 저관세이다.
\end{choices}
\end{question}
\answer{5}
